\chapter{Model Training}

\section{Notebook Scope}
Training is implemented in:
\begin{itemize}[leftmargin=2em]
\item \texttt{\detokenize{hand_gesture_model/train2.ipynb}}
\end{itemize}

This notebook builds a unified one-hand and two-hand classifier with fixed 126-feature input vectors.

\section{Preprocessing and Feature Extraction}
For each labeled bounding box:
\begin{enumerate}[leftmargin=2em]
\item Read YOLO box and crop hand region.
\item Detect up to two hands with MediaPipe.
\item Normalize each hand relative to wrist landmark.
\item Scale landmarks by the maximum 2D norm.
\item Assign features to left and right slots.
\item Zero-pad missing hand slot when only one hand is present.
\end{enumerate}

Normalization equation per hand:
\[
\hat{\mathbf{p}}_i = \mathbf{p}_i - \mathbf{p}_0,
\quad
s = \max_j \left\|\hat{\mathbf{p}}_{j,xy}\right\|_2,
\quad
\tilde{\mathbf{p}}_i = \frac{\hat{\mathbf{p}}_i}{\max(s,10^{-6})}
\]

Final feature vector:
\[
\mathbf{x} = [\tilde{\mathbf{p}}_{L,1..21},\tilde{\mathbf{p}}_{R,1..21}] \in \mathbb{R}^{126}
\]

\section{Training Configuration}
\begin{itemize}[leftmargin=2em]
\item Model: RandomForestClassifier
\item Trees: 400
\item Class balancing: \texttt{balanced\_subsample}
\item Split (v1, \texttt{train2.ipynb}): 80:20 stratified train-test over pooled samples
\item Split (v2, \texttt{train\_unified\_v2.py}): original Roboflow train/valid/test folders
\item Random seed: 42
\end{itemize}

\section{Evaluation Protocol v2 (Leak-Free Split)}
The v1 protocol pooled samples from all Roboflow folders and re-split them
randomly. Because the source dataset contains bursts of near-identical frames
of the same person performing the same sign, a random split can place sibling
frames on both sides of the train/test boundary and inflate the estimate.

The v2 protocol (\texttt{\detokenize{hand_gesture_model/train_unified_v2.py}})
instead honors the dataset's original split: 2{,}768 train, 798 validation,
and 395 test samples. Model selection uses validation macro-F1 only; the test
split is evaluated once for the final report.

\subsection{Landmark-Space Augmentation Study}
Two augmentation variants were evaluated on the training split only:
\begin{itemize}[leftmargin=2em]
\item \textbf{Mirror} (2$\times$): negate $x$ after wrist-centering and swap
the left/right hand blocks.
\item \textbf{Full} (4$\times$): mirror plus two rotations of $(x,y)$ around
the wrist by $\pm(5\text{--}15)^{\circ}$ with Gaussian jitter
($\sigma = 0.01$).
\end{itemize}
Neither variant improved validation macro-F1 beyond the selection margin
(0.005), and both \emph{reduced} held-out test performance (mirror:
96.46\%; full: 96.96\%; no augmentation: 97.47\%). The mirror result is
notable: hand signs are chirality-sensitive (which hand crosses on top is
class-relevant), so mirrored copies inject label noise. The deployed v2
model therefore uses the clean, non-augmented training set.

\section{Exported Artifacts}
Unified pipeline outputs are written to:
\begin{itemize}[leftmargin=2em]
\item \texttt{\detokenize{hand_gesture_model/naruto-hand-sign-4/phase1_outputs_unified/hand_landmarks_dataset_unified.csv}}
\item \texttt{\detokenize{hand_gesture_model/naruto-hand-sign-4/phase1_outputs_unified/X_landmarks_unified.npy}}
\item \texttt{\detokenize{hand_gesture_model/naruto-hand-sign-4/phase1_outputs_unified/y_labels_unified.npy}}
\item \texttt{\detokenize{hand_gesture_model/naruto-hand-sign-4/phase1_outputs_unified/confusion_matrix_unified.csv}}
\item \texttt{\detokenize{hand_gesture_model/naruto-hand-sign-4/phase1_outputs_unified/hand_gesture_rf_unified.mlmodel}}
\item \texttt{\detokenize{hand_gesture_model/naruto-hand-sign-4/phase1_outputs_unified/training_summary_unified.csv}}
\end{itemize}

The v2 pipeline additionally writes to
\texttt{\detokenize{phase1_outputs_unified_v2/}}:
\begin{itemize}[leftmargin=2em]
\item \texttt{\detokenize{hand_gesture_rf_unified.mlmodel}} (deployed to the iOS app)
\item \texttt{\detokenize{classification_report_v2.txt}} (per-class precision/recall/F1)
\item \texttt{\detokenize{confusion_matrix_v2.csv}}
\item \texttt{\detokenize{training_summary_v2.json}}
\end{itemize}

\section{Evaluation Metrics}
\begin{table}[H]
\centering
\caption{Unified model extraction and evaluation metrics (v2, leak-free split)}
\label{tab:evaluation_metrics}
\begin{tabular}{p{0.56\linewidth}p{0.34\linewidth}}
\toprule
Metric & Value \\
\midrule
Usable samples & 3961 \\
One-hand samples & 3289 (83.03\%) \\
Two-hand samples & 672 (16.97\%) \\
Skipped images & 2181 \\
Skipped no hands & 2177 \\
Feature dimension & 126 \\
Split & Roboflow train/valid/test (2768/798/395) \\
Test accuracy (v2, clean split) & \textbf{97.47\%} \\
Test macro-F1 (v2, clean split) & 0.9709 \\
Reference: v1 pooled random 80:20 accuracy & 96.85\% \\
CoreML export & Success (hand\_gesture\_rf\_unified.mlmodel) \\
\bottomrule
\end{tabular}
\end{table}


\section{Confusion Matrix Notes}
The exported confusion matrix shows strong diagonal dominance, with minor confusion among visually similar signs (for example, dragon vs snake, and ram vs tiger under some samples).

\placeholderfigure{confusion_matrix_placeholder}{figures/screenshots/confusion_matrix_heatmap.png}{Confusion matrix heatmap placeholder generated from CSV metrics.}
\placeholderfigure{training_notebook_outputs}{figures/screenshots/training_notebook_outputs.png}{Notebook output screenshot placeholder for extraction summary and training logs.}
